\chapter{Algebraic examples of ordered reconstruction}
\label{ch:algebraic-reconstruction-examples}

An augmentation selects a module whose derived endomorphisms can be computed by a bar construction.
A coproduct instead specifies how an algebra acts on tensor products of modules.
The examples below construct both kinds of structure and exhibit the relations that make them work.

Free and square-zero algebras give the first bar calculations. The quantum plane supplies a nontrivial Yoneda product.
Hall flags, rational exchange and reflection equations then provide explicit quantum-group coefficients.
Their realization by fields requires the field maps and collision compatibilities in addition to these algebraic formulas.

\section{Augmentations and two elementary bars}

Fix a field $k$ of characteristic zero.
An augmented algebra is a unital associative algebra $A$ with a unital algebra map
$\varepsilon:A\to k$. Write $I=\ker\varepsilon$.
For an ordinary algebra concentrated in degree zero, use the homological normalized bar
\[
 B_n(A)=I^{\otimes n},\qquad B_0(A)=k,
 \qquad
 d[a_1|\cdots|a_n]
 =\sum_{i=1}^{n-1}(-1)^{i-1}
 [a_1|\cdots|a_i a_{i+1}|\cdots|a_n].
\]

Deconcatenation cuts the tensor at every position, including the two empty ends.
Thus $B(A)$ is a coaugmented conilpotent differential graded coalgebra.
Conilpotence here means that repeated reduced coproducts kill each element.
In the cohomological convention of Chapter~\ref{mht:ch:signs-duals}, $B_n$ lies in degree $-n$.
All tensors and sums in this section are algebraic.

Let $V$ be finite-dimensional, give it internal weight one, and write $sV$ for $V$ in homological degree one.
The tensor algebra $T(V)=\bigoplus_{m\ge0}V^{\otimes m}$ is augmented by projection to $k$.
The square-zero algebra $k\oplus V$ has the same augmentation and product $V^2=0$.

\begin{proposition}[Free and square-zero bars]
\label{prop:alg-free-square-zero}
There is a quasi-isomorphism of coaugmented differential graded coalgebras
\[
 k\oplus sV\longrightarrow B(T(V)),\qquad sv\longmapsto[v],
\]
where $sV$ is primitive and its differential is zero.
For the square-zero algebra,
\[
 B(k\oplus V)=\bigoplus_{n\ge0}(sV)^{\otimes n}
\]
with zero differential and deconcatenation coproduct.
Consequently the first augmentation has square-zero positive Yoneda algebra $V^*$ in degree one. The second has tensor Yoneda algebra $T(V^*)$, with $V^*$ in degree one.
\end{proposition}
\begin{proof}
Fix an underlying tensor of $m$ factors of $V$.
A bar expression with that tensor is specified by a subset $C$ of its $m-1$ internal gaps.
The differential deletes the $i$th selected gap with sign $(-1)^{i-1}$.
For $m\ge2$, define $h(C)$ by inserting the first gap if it is absent, and by zero if it is present.

Deleting that newly inserted gap contributes $C$ to $dh(C)$. The other deletions cancel the terms of $hd(C)$.
If the first gap was already present, its deletion contributes $C$ to $hd(C)$.
Thus $dh+hd=1$ on every weight $m\ge2$.

Weights zero and one are $k$ and $sV$.
Their inclusion preserves the coproduct, proving the first assertion.

When $V^2=0$, every merged product in the bar differential vanishes.
Dualizing deconcatenation gives concatenation of dual tensors.
Dualizing the primitive coalgebra $k\oplus sV$ gives a square-zero positive ideal.
These dual models compute derived endomorphisms of the augmentation module and hence the Yoneda products.
\end{proof}

The grading determines the completion in these formulas.
The degree-$n$ Yoneda space of $k\oplus V$ is $(V^*)^{\otimes n}$. The ordinary graded algebra is their direct sum.
Its internal-weight completion is the different algebra
$\prod_{n\ge0}(V^*)^{\otimes n}$, with multiplication defined by the finite sum over two weights adding to a fixed weight.
A completion must therefore be specified even in this elementary example.

\section{Paths and the direction of extensions}

For a finite quiver $Q$, put $S=\bigoplus_{i\in Q_0}ke_i$, where the $e_i$ are orthogonal idempotents summing to one.
An arrow $a:i\to j$ is represented in the $S$-bimodule $M$ by $a=e_j a e_i$.
The algebra $A=T_S(M)$ consists of finite linear combinations of composable paths, multiplied with the rightmost path traversed first.
Projection to the paths of length zero is the augmentation $A\to S$.
The relative bar uses tensors over $S$.

\begin{proposition}[Relative path resolution]
\label{prop:alg-path-resolution}
The sequence of left $A$-modules
\[
 0\longrightarrow A\otimes_S M
 \xrightarrow{\ \mu\ }A\longrightarrow S\longrightarrow0,
 \qquad \mu(a\otimes m)=am,
\]
is a projective resolution.
The relative bar admits the coalgebra quasi-isomorphism
$S\oplus sM\to B_S(A)$ given by the paths of length zero and one.
\end{proposition}
\begin{proof}
Every nonempty path has a unique first arrow, which is its rightmost factor in the chosen multiplication convention.
Removing that factor gives the inverse of $\mu$ onto the augmentation ideal.
Semisimplicity of $S$ makes $M$ projective over $S$, so $A\otimes_S M$ is projective over $A$.
The gap contraction in Proposition~\ref{prop:alg-free-square-zero} respects endpoints and tensors over $S$.
It gives the relative coalgebra comparison.
\end{proof}

For $1\xrightarrow{a}2$, the projective left module $Ae_1$ has basis $e_1,a$ and radical $ka\cong S_2$.
The module $Ae_2=S_2$ is simple and projective.
Thus
\[
 0\longrightarrow Ae_2\xrightarrow{\ e_2\mapsto a\ }Ae_1
 \longrightarrow S_1\longrightarrow0
\]
gives $\Ext^1_A(S_1,S_2)=k$, while all other degree-one extensions between the two simples vanish.
Higher extensions vanish by the resolution.

There are two bimodule conventions to distinguish.
The vector-space dual $M^*=\Hom_k(M,k)$ with $(sft)(m)=f(tms)$ reverses endpoints:
$(e_jMe_i)^*=e_iM^*e_j$.
By contrast, in the algebra $E=\Ext^*_A(S,S)$ with product $\alpha\beta$ meaning first $\beta$ and then $\alpha$,
\[
 e_jE^1e_i=\Ext^1_A(S_i,S_j).
\]
The extension above therefore lies in $e_2E^1e_1$.
An identification with the endpoint-reversing dual must include the corresponding opposite-algebra convention.

\section{The quantum plane and its Yoneda product}

Let $q\in k^\times$ and set $A_q=k\langle x,y\rangle/(yx-qxy)$.
Assign $x,y$ internal weight one and use the augmentation killing them.
The rule $yx\mapsto qxy$ decreases the inversion number of a tensor.
Its leading pair has no overlap with itself, and reductions on disjoint pairs commute.
It follows that the monomials $x^ay^b$, $a,b\ge0$, form a basis.
In particular,
\[
 (x^ay^b)(x^cy^d)=q^{bc}x^{a+c}y^{b+d}.
\]

\begin{proposition}[Resolution and product]
\label{prop:alg-quantum-plane}
The augmentation has the linear free resolution
\[
 0\longrightarrow A_q e_r
 \xrightarrow{d_2}A_q e_x\oplus A_q e_y
 \xrightarrow{d_1}A_q\longrightarrow k\longrightarrow0,
\]
where $e_x,e_y$ have internal weight one, $e_r$ has weight two, and
\[
 d_1(e_x)=x,\qquad d_1(e_y)=y,\qquad
 d_2(e_r)=y e_x-qx e_y.
\]
With the composition convention above,
\[
 \Ext^*_{A_q}(k,k)
 =k\langle\xi,\eta\rangle/(\xi^2,\eta^2,\xi\eta+q\eta\xi),
 \qquad |\xi|=|\eta|=1.
\]
\end{proposition}
\begin{proof}
The relation gives $d_1d_2=0$, and $\operatorname{im}d_1$ is the augmentation ideal.
Write $u=\sum u_{ab}x^ay^b$ and $v=\sum v_{ab}x^ay^b$.
The equation $ux+vy=0$ is equivalent to
\[
 q^b u_{a-1,b}+v_{a,b-1}=0,
\]
with negative-index coefficients zero.
It gives $u_{a0}=0$ and $v_{ab}=-q^{b+1}u_{a-1,b+1}$.
For $c=\sum u_{a,b+1}x^ay^b$, these identities say $u=cy$ and $v=-qcx$.

This proves exactness at $A_qe_x\oplus A_qe_y$.
Right multiplication by $y$ is injective on the displayed basis, so $d_2$ is injective.

Applying $\Hom_{A_q}(-,k)$ kills the differential and gives dimensions $1,2,1$.
Represent the degree-one classes by maps to $k$ with
$\xi(e_x)=1$, $\xi(e_y)=0$, $\eta(e_x)=0$, $\eta(e_y)=1$.
Their degree-minus-one lifts to endomorphisms of the resolution satisfy
\[
 \xi(e_r)=-e_y,\qquad \eta(e_r)=q e_x.
\]
Indeed $d\xi+\xi d=d\eta+\eta d=0$ on $e_r$.
If $\theta(e_r)=1$ denotes the degree-two class, composition gives
\[
 \xi\eta=q\theta,\qquad \eta\xi=-\theta,
 \qquad \xi^2=\eta^2=0.
\]
These products give precisely the asserted four-dimensional algebra.
\end{proof}

The resolution is linear in internal weight, which is the Koszul property in this example.
Its dual relation can also be read directly from
$k(y\otimes x-q x\otimes y)$: the annihilator is spanned by
$\xi^2,\eta^2,\xi\eta+q\eta\xi$.
The explicit lifts fix the convention needed to identify that quadratic dual with the Yoneda algebra.

\section{Weyl and Clifford coefficients}

A scalar deformation of a quadratic relation can destroy the augmentation.
To see this without a completion ambiguity, first work over $R=k[\hbar]$ and form finite sums of tensor monomials.
For an even finite-dimensional space $V$ and alternating form $\omega$, define
\[
 W_\hbar(V,\omega)=T_R(V\otimes R)\big/
 (uv-vu-\hbar\omega(u,v)).
\]
For a purely odd space and symmetric $\omega$, define the Clifford algebra by
$uv+vu=\hbar\omega(u,v)$.
In particular, an odd generator satisfies $u^2=\hbar\omega(u,u)/2$.

\begin{proposition}[Ordered bases and scalar curvature]
\label{prop:alg-weyl-clifford}
After choosing an ordered basis of $V$, the Weyl algebra has an $R$-basis of nondecreasing monomials.
The Clifford algebra has the same basis with each generator occurring at most once.
For the filtration by the number of generators,
\[
 \operatorname{gr}W_\hbar=\operatorname{Sym}_R(V\otimes R),
 \qquad
 \operatorname{gr}\operatorname{Cl}_\hbar=\Lambda_R(V\otimes R).
\]
The quadratic relation space carries a scalar functional $c$ sending each relation to its prescribed value $\hbar\omega$.
If $c\ne0$, sending all generators to zero does not define an augmentation to $R$.
\end{proposition}
\begin{proof}
Replace each descending adjacent pair by its ordered pair and its scalar correction.
In the odd case also replace a repeated pair by its scalar square.
Order tensors first by length and then lexicographically.
Every replacement decreases this order.

In a triple of distinct generators, either reduction sequence makes the same three interchanges. Each scalar correction commutes with the remaining generator.
The resulting ordered terms and scalar terms therefore agree.

For repeated odd generators, the square rule agrees with the diagonal relation because $2u^2=\hbar\omega(u,u)$.
This resolves the remaining overlaps, proving the bases.

The scalar terms lower the filtration by two, so the associated graded algebras are as stated.
Applying a proposed augmentation to a relation would give $0=c(r)$, proving the last assertion.
\end{proof}

In a curved Koszul description, $c$ is the degree-two curvature element in the dual of the quadratic relation space. The presentation here has no linear correction and hence no linear dual differential.
The functional and its degree belong to this quadratic algebraic construction.
Identifying it with a geometric anomaly requires a map of the relevant deformation complexes.

One may subsequently complete in $\hbar$.
In that completion, an element can have unbounded tensor length as its $\hbar$-order increases.
The finite-length filtration above is not exhaustive on such elements.
The algebraic PBW basis describes this completion as formal series in $\hbar$.
Each coefficient is a finite linear combination of the stated monomials.

\section{Enveloping bars and their homology}

Let $\mathfrak g$ be an ordinary Lie algebra over $k$.
The universal enveloping algebra $U(\mathfrak g)$ has augmentation killing $\mathfrak g$.
Its normalized bar is generally much larger than its Chevalley chain complex
$C_n(\mathfrak g)=\Lambda^n\mathfrak g$.
Set
\[
 d(x_1\wedge\cdots\wedge x_n)
 =\sum_{i<j}(-1)^{i+j+1}[x_i,x_j]\wedge
 x_1\wedge\cdots\widehat{x_i}\cdots\widehat{x_j}\cdots\wedge x_n.
\]
The exterior coalgebra coproduct splits the factors into two ordered sublists with the permutation sign.

\begin{theorem}[Ordinary bar--Chevalley comparison]
\label{thm:alg-bar-chevalley}
Antisymmetrization defines a coalgebra quasi-isomorphism
\[
 \kappa:C_*(\mathfrak g)\longrightarrow B(U(\mathfrak g)),\qquad
 \kappa(x_1\wedge\cdots\wedge x_n)
 =\sum_{\sigma\in\Sigma_n}\operatorname{sgn}(\sigma)
 [x_{\sigma(1)}|\cdots|x_{\sigma(n)}].
\]
If $\mathfrak g$ is positively graded and its bracket adds weights,
this comparison is a quasi-isomorphism in each weight and on the degreewise product over weights.
The product completion uses tensors completed by total weight.
\end{theorem}
\begin{proof}
Pair terms whose merged adjacent entries occur in opposite orders.
Their difference is the Lie bracket, with the sign in the displayed Chevalley differential.
This proves the chain-map equation.
Cutting an antisymmetrized tensor is the same signed sum as choosing its left subset and antisymmetrizing both parts.
Hence $\kappa$ preserves coproducts.

Filter by total PBW degree, putting $\Lambda^n\mathfrak g$ in filtration degree $n$.
The associated graded comparison is
$\Lambda(V)\to B(\operatorname{Sym}(V))$, for $V=\mathfrak g$.
It is the comparison of the Koszul and bar free resolutions of $k$.
For completeness, on polynomial degree $p$ and exterior degree $q$ with $p+q>0$, the Koszul contraction is
\[
 h=\frac1{p+q}\sum_i (x_i\wedge -)\frac{\partial}{\partial x_i},
 \qquad \delta=\sum_i x_i\iota_i,
 \qquad \delta h+h\delta=1.
\]
Here $\iota_i$ contracts the exterior generator $x_i$. Only finitely many summands act on each element.

The free bar resolution is contracted by its extra degeneracy.
The augmentation-preserving comparison between the two free resolutions has a homotopy inverse obtained by lifting basis elements degree by degree.
Tensoring those homotopies with $k$ proves the associated graded assertion.
Multiplication by $(-1)^n$ in degree $n$ identifies the two overall bar-differential sign conventions.

A cycle in the cone of $\kappa$ has finite PBW degree.
A primitive for its highest associated graded term lowers that degree.
Iteration terminates, proving acyclicity of the cone.
In positive total weight $w$, PBW degree is at most $w$, so the argument also works weight by weight.
Products of complexes of vector spaces have componentwise kernels and images.
The comparison consequently extends to the stated product completion, which is conilpotent in each weight.
\end{proof}

For every augmented algebra, $H_1B(A)=I/I^2$.
For an enveloping algebra this becomes
\[
 H_1B(U(\mathfrak g))=\mathfrak g/[\mathfrak g,\mathfrak g].
\]
The inclusion of generators is surjective onto $I/I^2$ and kills brackets.
Its inverse follows from the augmented algebra map
$U(\mathfrak g)\to k\oplus\mathfrak g/[\mathfrak g,\mathfrak g]$ with square-zero second summand.
This degree-one calculation often decides a proposed homology formula before any higher calculation is needed.

For $\mathfrak{sl}_2=\langle e,h,f\rangle$, with
$[h,e]=2e$, $[h,f]=-2f$ and $[e,f]=h$, the differential
$\Lambda^2\mathfrak{sl}_2\to\mathfrak{sl}_2$ sends
\[
 e\wedge h\mapsto-2e,\qquad e\wedge f\mapsto h,
 \qquad h\wedge f\mapsto-2f.
\]
It is an isomorphism, while
$d(e\wedge h\wedge f)=-2e\wedge f-2f\wedge e=0$.
The homology polynomial is therefore $1+t^3$.
The Chevalley chain dimensions are $1,3,3,1$.

The bar chain space $B_n(U(\mathfrak{sl}_2))$, however, is countably infinite-dimensional for every $n\ge1$.
Its finite PBW pieces have generating function
\[
 \sum_{d\ge0}\dim\operatorname{gr}^{\mathrm{PBW}}_dB_n(U(\mathfrak{sl}_2))z^d
 =\bigl((1-z)^{-3}-1\bigr)^n.
\]
Thus $3^n$ counts only the least PBW degree $d=n$.

Over $\mathbb C$, the cohomology of a finite-dimensional simple Lie algebra of rank $r$ is exterior on generators in degrees $2m_i+1$, where $m_i$ are its exponents.
This standard Lie-cohomology theorem gives
\[
 P_{\mathfrak{sl}_3}(t)=(1+t^3)(1+t^5),\qquad
 P_{\mathfrak g_2}(t)=(1+t^3)(1+t^{11}).
\]
See Etingof's \emph{Lie groups and Lie algebras, I}, \S44.
The invariant-form comparison, rank-two examples and exponent formula appear in Theorem~44.12, Example~44.16 and Theorem~44.17.
The equality of homology and cohomology dimensions here uses finite-dimensional Chevalley chain spaces.

For the positive Witt algebra
\[
 L_1=\bigoplus_{m\ge1}\mathbb C e_m,\qquad [e_i,e_j]=(j-i)e_{i+j},
\]
the abelianization has basis $e_1,e_2$.
Indeed brackets have weight at least three and
$e_m=(m-2)^{-1}[e_1,e_{m-1}]$ for $m\ge3$.

The higher calculation is Goncharova's theorem: for every $n\ge1$, the homology of $B(U(L_1))$ has one class in each weight
\[
 \frac{3n^2-n}{2},\qquad \frac{3n^2+n}{2},
\]
and vanishes in other weights.\footnote{L.~V.~Goncharova, \emph{The cohomologies of Lie algebras of formal vector fields on the line}. Functional Analysis and Its Applications \textbf{7} (1973), 91--97 and 194--203. The weight formula is also stated as Theorem~1.4 of D.~Millionschikov, \emph{Algebra of formal vector fields on the line and Buchstaber's conjecture}, arXiv:0807.0653.}

Weightwise duality and Theorem~\ref{thm:alg-bar-chevalley} transfer the Lie-cohomology formula to this bar.
Its Euler series is consequently
\[
 \prod_{m\ge1}(1-z^m)
 =1+\sum_{n\ge1}(-1)^n
 \bigl(z^{(3n^2-n)/2}+z^{(3n^2+n)/2}\bigr).
\]
The homology theorem supplies the degree and multiplicity of each class. The Euler series alone does not determine them.

\section{Planar path counts and bar gradings}

A Motzkin path of length $n$ stays at nonnegative height, starts and ends at height zero, and has steps $(1,1),(1,0),(1,-1)$.
Let $M(n)$ count these paths and $R(n)$ count those with no horizontal step at height zero.
Both include the empty path.
Write $M(z)=\sum M(n)z^n$ and $R(z)=\sum R(n)z^n$ in $\mathbb Q[[z]]$.

A first-return decomposition gives
\[
 M=1+zM+z^2M^2,\qquad R=1+z^2MR.
\]
In the second decomposition, the interior of the first excursion is arbitrary, while the remainder obeys the height-zero restriction.
The quadratic equation implies
$(1+zM)(1-z^2M)=1+z$ and therefore
\[
 (1+z)R=1+zM,\qquad
 M(n)=R(n)+R(n+1).
\]

Solving with constant term one gives
\[
 M(z)=\frac{1-z-\sqrt{1-2z-3z^2}}{2z^2},
 \qquad R(z)=\frac{1+zM(z)}{1+z}.
\]
The square root has constant term one, and the apparent pole is removable.

For $n\ge1$, the derived sequences
\[
 a_n=M(n+1)-M(n)=1,2,5,12,30,76,\ldots,
 \qquad b_n=R(n+3)=3,6,15,36,91,\ldots
\]
belong to the same quadratic generating-function field.
They nevertheless disagree already in degree one with the proposed enveloping homologies:
\[
 \dim H_1B(U(L_1))=2\ne a_1,
 \qquad \dim H_1B(U(\mathfrak{sl}_2))=0\ne b_1.
\]
They also cannot count the full bar chain spaces, whose positive degrees are countably infinite-dimensional.

For $L_1$, the finite weightwise chain count is instead
\[
 \sum_{w\ge0}\dim B_n(U(L_1))_w z^w
 =\left(\prod_{m\ge1}(1-z^m)^{-1}-1\right)^n.
\]
A comparison with a path complex must specify its differential, grading and map to this bar.
The choice of $U(L_1)$ is itself algebraic data. A Virasoro coefficient or central charge does not determine a transverse multiplication.

\section{Hall multiplication for one arrow}

Let $Q:1\to2$ over $\mathbb F_q$, and choose $v$ with $v^2=q$ in a characteristic-zero coefficient field.
A representation is a linear map between its two vertex spaces.
For isomorphism classes $[A],[B]$, define twisted Hall multiplication by
\[
 [A]*[B]=v^{\langle\dim A,\dim B\rangle}
 \sum_{[L]}F^L_{AB}[L],\qquad
 \langle(a,b),(c,d)\rangle=ac+bd-ad,
\]
where $F^L_{AB}$ counts subrepresentations isomorphic to $B$ with quotient isomorphic to $A$.
Associativity follows by counting the same two-step flags in either order. Bilinearity of the Euler form makes the two total twisting exponents equal.

Write $S_1,S_2$ for the vertex simples, $D=S_1\oplus S_2$, and $M$ for the nonzero map $\mathbb F_q\to\mathbb F_q$.
The two representations of dimension $(1,1)$ give
\[
 E_1E_2=v^{-1}([D]+[M]),\qquad E_2E_1=[D],
 \qquad E_i=[S_i].
\]
In particular $E_1E_2-v^{-1}E_2E_1=v^{-1}[M]$.

\begin{proposition}[The two quantum Serre relations]
\label{prop:alg-hall-serre}
The Hall generators satisfy
\[
 E_1^2E_2-(v+v^{-1})E_1E_2E_1+E_2E_1^2=0
\]
and the relation obtained by exchanging $1$ and $2$.
\end{proposition}
\begin{proof}
A representation of dimension $(2,1)$ is a map $f:\mathbb F_q^2\to\mathbb F_q$, of rank zero or one.
Read the successive simple quotients of a complete flag from top to bottom.
The counts are
\[
\begin{array}{c|ccc}
 \operatorname{rank}f&112&121&211\\\hline
 0&q+1&q+1&q+1\\
 1&q+1&1&0.
\end{array}
\]
In rank zero, each of the $q+1$ lines is allowed.
In rank one, the middle pattern forces the kernel line, the last pattern would force $f=0$, and the first allows every line.

The Euler exponents of these three patterns are $-1,0,1$.
The Serre coefficient on the rank-zero representation is
$(q+1)(v^{-1}-(v+v^{-1})+v)=0$.
On the rank-one representation it is
$v^{-1}(q+1)-(v+v^{-1})=0$.

Dualizing representations reverses flags and the arrow. Interchanging the vertices restores $Q$.
The resulting anti-isomorphism exchanges $E_1,E_2$ and sends this relation to the other one.
\end{proof}

The relation is an exact count of extensions with its stated normalization.
Identifying a Hall algebra of a geometric category with these coefficients additionally requires an extension-preserving comparison of categories and of the counting or integration measures.

\section{Exchange in an ordered chamber}

Fix a finite set of pairwise distinct spectral labels and a total order on it.
Let $K$ be a coefficient field containing their differences, and let each label carry a copy of a finite-dimensional space $V$.
An exchange rule replaces an inverted adjacent pair by a $K$-linear combination with the two labels in increasing order.
Write its coefficient as an invertible matrix $R(u-v)$. Tensor positions track the original labeled factors.
Apply rules only to inversions.

\begin{proposition}[Chamber confluence]
\label{prop:alg-exchange}
Suppose no further relations are imposed on the ordered tensor spaces.
The exchange rules are confluent precisely when, on each triple of distinct labels,
\[
 R_{12}(u-v)R_{13}(u-w)R_{23}(v-w)
 =R_{23}(v-w)R_{13}(u-w)R_{12}(u-v).
\]
When this equation holds, the ordered tensors form a basis of the quotient.
\end{proposition}
\begin{proof}
Every exchange decreases the inversion number of the label sequence.
Disjoint exchanges commute.
For three descending labels, the two sequences of three exchanges have the displayed coefficients.
These are all overlapping ambiguities, so equality resolves every reduction by induction on the inversion number.
Conversely, confluence and independence of the ordered tensor space force equality of the two coefficient operators.
\end{proof}

The rational permutation matrix supplies the basic example.
Let $P(v\otimes w)=w\otimes v$ and set
\[
 R(u)=1-\frac{\hbar P}{u}.
\]
The Yang--Baxter identity follows in the group algebra of $\Sigma_3$.
The terms of degree zero and one in $\hbar$ agree directly.
At degree two, the two three-cycles have coefficients equal by
\[
 \frac1{u(u+v)}+\frac1{v(u+v)}=\frac1{uv};
\]
at degree three the equality is
$P_{12}P_{13}P_{23}=P_{23}P_{13}P_{12}$.
This proves the operator identity on $V^{\otimes3}$ in every dimension.

If exchange is also imposed in the reverse direction, its compatibility requires
$R_{21}(-u)R(u)=1$ on the chosen pair space.
For the rational matrix,
\[
 R_{21}(-u)R(u)=\left(1-\frac{\hbar^2}{u^2}\right)1.
\]
Over $k(\hbar,u)$, imposing both relations therefore kills every degree-two product between those distinct labels.
Even $\dim V=1$ exhibits this failure, although the Yang--Baxter equation is then automatic.
Over a formal coefficient ring, it instead produces $\hbar^2$-torsion.

The scalar normalization
\[
 \widetilde R(u)=f(u)R(u),\qquad
 f(u)=\left(1-\frac{\hbar^2}{u^2}\right)^{-1/2}
 \in k(u)[[\hbar]]
\]
restores the reverse-crossing identity.
The series $f$ is even in $u$ and has constant term one. Its scalar nature preserves Yang--Baxter.
Its formal coefficient ring is part of the normalization.

\section{RTT multiplication, coproduct and PBW}

Work first over $R=k[\hbar]$.
Define $Y_\hbar(\mathfrak{gl}_N)$ by generators $t_{ij}^{(r)}$, $r\ge1$, and put
\[
 T(u)=I+\sum_{r\ge1}T^{(r)}u^{-r},\qquad
 t_{ij}^{(0)}=\delta_{ij}.
\]
The defining RTT relation is
\[
 R(u-v)T_1(u)T_2(v)=T_2(v)T_1(u)R(u-v),
 \qquad R(u)=I-\frac{\hbar P}{u}.
\]
It means the coefficient identities after clearing $u-v$:
\begin{equation}
\label{eq:alg-rtt-coefficients}
 [t_{ij}^{(r+1)},t_{kl}^{(s)}]-[t_{ij}^{(r)},t_{kl}^{(s+1)}]
 =\hbar\bigl(t_{kj}^{(r)}t_{il}^{(s)}-t_{kj}^{(s)}t_{il}^{(r)}\bigr),
 \qquad r,s\ge0.
\end{equation}
Indeed these are the coefficients of $u^{-r}v^{-s}$ in
$(u-v)[t_{ij}(u),t_{kl}(v)]$ and its numerator.
Each coefficient is a finite algebraic expression.

\begin{proposition}[Hopf structure and evaluation]
\label{prop:alg-rtt-hopf}
There is an algebraic Hopf structure with
\[
 \Delta(t_{ij}^{(r)})=
 \sum_{a=1}^N\sum_{p=0}^r t_{ia}^{(p)}\otimes t_{aj}^{(r-p)},
 \qquad \varepsilon(T(u))=I,\qquad S(T(u))=T(u)^{-1}.
\]
Let $U_\hbar(\mathfrak{gl}_N)$ be generated by $E_{ij}$ with
$[E_{ij},E_{kl}]=\hbar(\delta_{kj}E_{il}-\delta_{il}E_{kj})$.
For a central parameter $a$, evaluation is the homomorphism
\[
 T(u)\longmapsto I+\frac{E}{u-a},
 \qquad t_{ij}^{(r)}\longmapsto a^{r-1}E_{ij}.
\]
The inclusion $E_{ij}\mapsto t_{ij}^{(1)}$ has evaluation at $a=0$ as a left inverse.
\end{proposition}
\begin{proof}
If matrices $T,U$ satisfy RTT and their coefficient algebras commute, then
\[
 R(T_1U_1)(T_2U_2)
 =RT_1T_2U_1U_2
 =T_2T_1RU_1U_2
 =T_2U_2T_1U_1R.
\]
Here only factors from distinct coefficient algebras and distinct auxiliary positions have been interchanged.
Thus their matrix product satisfies RTT.
This proves the coproduct identity. Associativity of matrix multiplication proves coassociativity.
The counit follows by substituting $I$.

Since $T(u)=I+O(u^{-1})$, its inverse coefficient of order $r$ is a finite sum of products of coefficients of total order $r$.
Inverting the RTT identity shows that this inverse satisfies the relations for an antihomomorphism.
The two matrix inverse identities are the antipode axioms.
All these formulas lie in the ordinary coefficient algebra and tensor product. Only their generating series use $u^{-1}$-adic notation.

Substituting $I+E/(u-a)$ into RTT reduces it to the displayed commutator of the $E_{ij}$.
Equation~\eqref{eq:alg-rtt-coefficients} with $r=0,s=1$ gives that commutator for the first modes.
Evaluation fixes those modes, proving the claimed injection.
\end{proof}

\begin{theorem}[PBW for the rational RTT algebra]
\label{thm:alg-rtt-pbw}
For any total order on the generators, their ordered monomials form an $R$-basis of $Y_\hbar(\mathfrak{gl}_N)$.
For the filtration $\deg t_{ij}^{(r)}=r-1$,
\[
 \operatorname{gr}Y_\hbar(\mathfrak{gl}_N)
 \cong U_\hbar(\mathfrak{gl}_N[t]).
\]
For the different filtration $\deg t_{ij}^{(r)}=r$, the associated graded algebra is the polynomial algebra on the same generators.
\end{theorem}
\begin{proof}
Iterating the coefficient recurrence gives
\[
 [t_{ij}^{(r)},t_{kl}^{(s)}]
 =\hbar\sum_{b=1}^{\min(r,s)}
 \left(t_{kj}^{(b-1)}t_{il}^{(r+s-b)}
       -t_{kj}^{(r+s-b)}t_{il}^{(b-1)}\right).
\]
In the degree-$r$ filtration, each term on the right has strictly smaller total degree than the commutator on the left.
Sorting therefore proves spanning by ordered monomials.
In the degree-$r-1$ filtration, the $b=1$ term is the current Lie bracket and all other terms have lower degree.
There is consequently a surjection
$U_\hbar(\mathfrak{gl}_N[t])\to\operatorname{gr}Y_\hbar$.

To prove injectivity, iterate the coproduct and evaluation at independent parameters $a_1,\ldots,a_m$.
The resulting matrix is
\[
 \prod_{j=1}^m\left(I+\frac{E^{[j]}}{u-a_j}\right),
\]
with the indicated order and commuting coefficient algebras in distinct factors.
Give the parameters degree one and all $E^{[j]}$ degree zero.
The highest parameter-degree part of the image of $t_{ij}^{(r)}$ is
$\sum_{b=1}^m a_b^{r-1}E_{ij}^{[b]}$. A product using at least two nonidentity matrix factors has degree at most $r-2$.
Thus the induced map on current enveloping algebras is evaluation at these $m$ independent parameters.

These evaluation maps, as $m$ varies, are jointly injective.
Indeed, take a nonzero element of highest PBW degree $d$.
For $m=d$, pass to the PBW associated graded in each tensor factor and take the component of degree one in every factor.
A symmetric tensor of $d$ currents maps to its polarization, with the power of $a_b$ recording the current exponent in factor $b$.
Distinct basis tensors have distinct coefficients of these independent parameter monomials.

Polarization is injective in characteristic zero.
PBW over $R$ applies because the current Lie algebra is free over the $\mathbb Q$-algebra $R$.
Hence the original element has nonzero image in one of these evaluations.

The current enveloping surjection is injective, proving the first associated graded formula.
Its PBW basis proves independence of the ordered Yangian monomials and hence also the second formula.
\end{proof}

For $k=\mathbb C$, the usual complex Yangian is obtained by setting $\hbar=1$.
In an evaluation into the ordinary $U(\mathfrak{gl}_N)[\hbar]$, the matrix is instead $I+\hbar e/(u-a)$, where $[e_{ij},e_{kl}]=\delta_{kj}e_{il}-\delta_{il}e_{kj}$.
These conventions keep the deformation parameter in its proper location.
Molev's \emph{Yangians and their applications}, \S\S2.1--2.5, gives the standard RTT, evaluation, PBW and Hopf formulations.\footnote{A.~I.~Molev, \emph{Yangians and their applications}, Handbook of Algebra \textbf{3} (2003), 907--959. arXiv:math/0211288, Proposition~2.2 and Theorems~2.6, 2.8.}

\section{An explicit reflection algebra}

A reflection equation is a different quadratic relation.
For a constant invertible $R\in\operatorname{End}(V^{\otimes2})$, put $R_{21}=PRP$ and define its coefficient algebra by
\[
 R_{21}K_1R K_2=K_2R_{21}K_1R.
\]

Even for Yang--Baxter solutions, this equation does not determine a universal ordered-monomial theorem.
For $R=I$, it says that all entries of $K$ commute.
For $R=P$, it is the identity $K_2^2=K_2^2$ and imposes no relations at all.
If $\dim V>1$, the resulting symmetric and free algebras have different Hilbert series.
A PBW assertion must therefore specify the matrix and the rewriting rules.

Take $V=k^2$, $q\in k^\times$, and use the basis $11,12,21,22$ of $V^{\otimes2}$.
Set
\[
 R=\begin{pmatrix}
 q&0&0&0\\
 0&1&q-q^{-1}&0\\
 0&0&1&0\\
 0&0&0&q
 \end{pmatrix},\qquad
 K=\begin{pmatrix}a&b\\c&d\end{pmatrix},\qquad \gamma=1-q^{-2}.
\]
This $R$ is invertible and satisfies the constant Yang--Baxter equation.

\begin{proposition}[Reflection PBW basis]
\label{prop:alg-reflection-pbw}
The reflection equation for this matrix is equivalent to
\begin{align*}
 ba&=q^2ab,&ca&=q^{-2}ac,&da&=ad,\\
 cb&=bc+\gamma(a^2-ad),&
 db&=bd+\gamma ab,&
 dc&=cd-\gamma q^{-2}ac.
\end{align*}
The monomials $a^i b^j c^k d^\ell$ form a basis, for all $q\ne0$.
\end{proposition}
\begin{proof}
Multiplying the four-by-four matrices gives the six displayed relations. Each remaining entry of the reflection equation is a scalar multiple of one of them or reduces by $da=ad$.
Conversely, the six relations make every matrix entry vanish.

Orient them from left to right with $a<b<c<d$.
Each rule decreases degree-lexicographic order.
The only overlapping leading pairs occur in $cba,dba,dca,dcb$.
Both reduction sequences give, respectively,
\[
\begin{array}{c|l}
 cba&abc+\gamma a^3-\gamma a^2d\\
 dba&q^2abd+q^2\gamma a^2b\\
 dca&q^{-2}acd-\gamma q^{-4}a^2c\\
 dcb&bcd+\gamma a^2d-\gamma ad^2.
\end{array}
\]

Disjoint replacements commute.
Induction on the decreasing order therefore gives a unique normal expression in every degree, proving the basis.
The constant Yang--Baxter identity for the displayed $R$ is verified by multiplication on $(k^2)^{\otimes3}$. It is independent of this explicit overlap argument.
\end{proof}

The three constructions now have specified domains.
Chamber exchange acts on labeled tensor spaces.
RTT gives an algebra and its tensor-product action through a coproduct.
Reflection gives the stated boundary coefficient algebra.
A physical identification must supply representations, supported field operations, boundary conditions and compatible collision maps in addition to these coefficients.
